\begin{abstract}

% Advanced multimodal AI agents can now collaborate with users to solve challenges in the world.  We explore eye tracking's role in such interaction to convey a user's attention relative to the physical environment.  We hypothesize that this knowledge improves contextual understanding for AI agents.  By observing hours of human-object interactions, we first measure the relationship between an eye tracker's signal quality and its ability to reliably place gaze on nearby physical objects.  We then conduct experiments which relay the user's scanpath history as additional context querying multimodal agents.  Our results show that eye tracking provides high value as a user attention signal and can convey information about the user's current task and interests to the agent.

Advanced multimodal AI agents can now collaborate with users to solve challenges in the world.  Yet, these emerging contextual AI systems rely on explicit communication channels between the user and system.  We hypothesize that implicit communication of the user's interests and intent would reduce friction and improve user experience when collaborating with AI agents.  In this work, we explore the potential of wearable eye tracking to convey signals about user attention.  We measure the eye tracking signal quality requirements to effectively map gaze traces to physical objects, then conduct experiments that provide visual scanpath history as additional context when querying vision language models.  Our results show that eye tracking provides high value as a user attention signal and can convey important context about the user's current task and interests, improving understanding of contextual AI agents.

\end{abstract}